\chapter{Preface}

This book presents reinforcement learning through the relations that define the problem, the derivations of those relations, the algorithms that apply them, and short Python programs that follow the same steps. It is intended for readers who know basic probability and statistics and who can read Python. A previous course in reinforcement learning is not required.

Each method begins with the object to be learned. The compact equation is stated next. When that equation abbreviates a longer expansion, the expansion is derived in the chapter or in an appendix from the probability identities in Appendix A, which are cited by number. The algorithm then records the update as a loop. A Python listing on a small environment follows that loop in the same section.

Appendix A contains the facts used later, including total probability, expectation, variance, importance sampling, and the Markov property. Those facts may be read first, or Chapter 1 may be started at once and the appendix consulted when an identity is cited. Longer derivations that would interrupt a chapter are placed in later appendices. They include the Bellman expansions, the Robbins--Monro step-size conditions, the policy-gradient theorem, and the non-negativity of KL, and they use the same identities as Appendix A.

The chapters treat Markov decision processes and the Bellman equations, dynamic programming, Monte Carlo methods, temporal-difference learning and control, function approximation, deep Q-networks, policy gradients, and actor--critic methods, and then reinforcement learning from human feedback together with InstructGPT-style and DeepSeek-R1-style training. A summary collects the objects. The algorithms, theorems, and named systems are those of the works listed in the References. The text is not a research monograph.

Each algorithm in the main chapters is accompanied by a program on a small state set. The listing stands next to the boxed steps so that one update can be read as a program. Two exercises at the end of the book apply several methods to one environment. The first is a $4\times 4$ GridWorld under value iteration, Monte Carlo control, SARSA, and Q-learning. The second is CartPole with a deep Q-network. Solutions are given. Chapters 11 and 12 use a language-model example with checkable answers. They describe the shape of InstructGPT-style and DeepSeek-R1-style training and do not train the published systems.
